
%\clearpage
\section{METODOLOGIA}

A organização metodológica do \NOMECURSO observa as Diretrizes Curriculares Nacionais para os cursos de Engenharia, definidas pela Resolução CNE/CES nº 2, de 24 de abril de 2019, e complementadas pela Resolução CNE/CES nº 1, de 26 de março de 2021. Atende, ainda, às orientações da Resolução CONSUP/IFCE nº 259, de 08 de janeiro de 2025, que estabelece o Manual de Normalização dos Projetos Pedagógicos dos Cursos de Engenharia do IFCE, bem como às diretrizes do SINAES (2017) para organização e avaliação dos cursos de graduação.

Os procedimentos metodológicos adotados nas unidades curriculares são descritos nos Programas de Unidade Didática (PUDs) e articulam métodos tradicionais consolidados no ensino de engenharia, tais como aulas expositivas e resolução estruturada de exercícios, com práticas orientadas ao desenvolvimento de competências, incluindo trabalhos dirigidos, estudos de caso, atividades experimentais em laboratórios, uso de ambientes computacionais e tecnologias digitais.

As metodologias contemplam, obrigatoriamente, atividades práticas e experimentais, em conformidade com o núcleo de formação proposto pelas DCNs da Engenharia, buscando garantir a articulação entre fundamentos científicos, domínio tecnológico e aplicação profissional. A integração teoria–prática é reforçada por projetos, estudos aplicados e vivências técnicas no campus ou em ambientes produtivos externos, incluindo visitas técnicas, estágios e participação em grupos de pesquisa.

A interdisciplinaridade é princípio estruturante da formação. Cada unidade curricular identifica, em seu PUD, os vínculos conceituais e procedimentais com demais componentes da matriz, orientando o estudante a compreender a engenharia como campo integrado e evolutivo. O currículo também permite a cursabilidade de disciplinas em outros cursos do IFCE ou em áreas correlatas, ampliando a visão sistêmica e a capacidade de atuação em contextos multidisciplinares.

Conforme o Guia de Curricularização da Extensão do IFCE, o curso assegura, no mínimo, 10\% da carga horária total em atividades de extensão, integradas ao território e às demandas sociais, desenvolvidas por meio de programas, projetos, cursos, oficinas, eventos ou prestação de serviços tecnológicos. Tais atividades são planejadas para promover interação transformadora entre instituição e sociedade, articulando formação acadêmica, responsabilidade social e desenvolvimento local e regional.

Os conteúdos relativos à educação ambiental (Lei nº 9.795/1999), educação em direitos humanos (Resolução CNE/CP nº 1/2012) e relações étnico-raciais (Resolução CNE/CP nº 1/2004) são contemplados de forma transversal e também em componentes curriculares específicos, tais como Ética e Cidadania e Projetos Sociais, podendo ser aprofundados por meio de projetos de extensão, pesquisa aplicada e eventos acadêmicos.

Nas disciplinas com atividades de laboratório, o estudante interage diretamente com equipamentos, instrumentação, \textit{softwares} de simulação e sistemas automatizados, realizando coleta e análise de dados, ensaios e elaboração de relatórios técnicos, promovendo o desenvolvimento de autonomia, rigor metodológico e precisão técnica.

A utilização das Tecnologias da Informação e Comunicação (TICs) integra-se ao processo formativo por meio de ambientes virtuais de aprendizagem, recursos de videoconferência, bibliotecas digitais, plataformas de simulação e ferramentas colaborativas, assegurando acessibilidade, flexibilidade e continuidade dos estudos.

As ações de monitoria, orientadas por docentes, reforçam o acompanhamento acadêmico, contribuindo para o desenvolvimento de estratégias de estudo e a consolidação de conteúdos. O atendimento a estudantes com necessidades específicas segue os dispositivos legais vigentes, com acompanhamento institucional por equipe multiprofissional.

Outras ações de apoio ao discente estão descritas no item \ref{APOIO}, \textit{Apoio ao Discente}.
